% 【本文件写什么】impl 实现层：dummy
% - 定位：信号 impl 占位，未挂聚合 feature 链。
% - crate 路径：os/components/wateros-ipc/ipc-signal/signal-impl/impl-dummy/src/
% - 如何实现对应 api-v0 trait：关键类型、状态机、数据结构
% - 算法或硬件交互流程，可用伪代码/流程图
% - 本 impl 启用的 Cargo [features] 与 arch feature，impl-riscv64 等
% - 与其它 impl 变体的差异、切换 feature 时的影响
% - 性能/内存假设与已知限制
% 【事实来源】
% - 上述 crate 源码与 Cargo.toml
% - os/feature-tree.txt 中指向本 impl 的 feature 链

\subsubsection{impl-dummy}

\code{wateros-ipc-signal-impl-dummy}为链接桩 crate，仅导出\code{add}占位函数，不实现\code{SignalRegistry}、不投递信号、不操纵线程掩码。

\begin{rustcode}
/// 占位算术：无信号投递语义。
#[inline]
pub fn add(left : u64, right : u64) -> u64 { left + right }
\end{rustcode}

根\code{wateros-ipc}的\code{feature signal}依赖\code{ipc-signal}聚合 crate 及其\code{signal-api/api-v0}，未将\code{impl-dummy}选为\code{active\_impl}；生产路径使用聚合\code{lib.rs}内的\code{SignalRegistry}。dummy 保留在 workspace 供未来拆分 impl 变体或最小 feature 组合编译。
